\documentclass[11pt]{article}
\usepackage{anysize,amsmath,url}
\usepackage{HW}
\marginsize{2.2cm}{2.2cm}{1.8cm}{1.8cm}
\begin{document}
\titleline{\semester}
\exhead{Task 02 -- Reflection Essay}

\noindent
Going into this task I assumed the deep-learning part would be the
mountain and the geometry a footnote. It turned out the other way round.
Running YOLOv8 on the avenue photo and getting a labelled bicycle box
back was a single function call. What actually took thought was believing
-- and then proving to myself -- that a flat photograph with no depth
information can still yield a real distance in metres.

The hinge of the whole task is the cross ratio. Perspective destroys
ordinary measurement: the avenue's two edges visibly converge towards the
arch, so equal real lengths shrink as they recede, and ``measure pixels,
multiply by a constant'' simply does not work. The cross ratio is the one
combination of four collinear points that a perspective camera leaves
unchanged, so its value in the photo equals its value on the ground. That
single invariant is what lets three known distances pin down a fourth
unknown one. Seeing the link back to the 25-May surveillance-camera
example -- locating a train between two stations using the point where the
rails meet -- was the moment the idea clicked.

Two practical lessons stuck with me. First, the arch and the vanishing
point are essentially the same point in this image, so I could not use
both as references; their pixel positions were too close and the
arithmetic went unstable. Moving to the intersection as the middle
reference cured it. Second, and more striking, the choice of point inside
the YOLO box mattered enormously. The box is a rectangle but the bicycle
only meets the ground at the bottom; using the box centre instead of the
wheel contact moved the answer by more than twenty metres. Near the
vanishing point one pixel is worth roughly $0.16$~m, so where you ``stand''
the object is not a detail -- it is the dominant source of error. By
contrast, rounding my pixel readings to the nearest ten changed the
result by under a metre. The modelling assumption, not the measurement
precision, was the thing to get right.

I cross-checked the bicycle distance two independent ways, once treating
the vanishing point as the point at infinity and once with three finite
landmarks, and the two agreed to within a few metres. That agreement did
more for my confidence than any single calculation could have.

If I were to redo it, I would spend longer nailing the vanishing point,
since real path edges are not perfectly straight and that uncertainty
propagates into everything downstream. The broader takeaway is that the
machine-learning tool is a convenience that locates an object; the
understanding lives entirely in the projective geometry and in respecting
the assumptions -- collinearity, a level camera, an object truly on the
line -- that the cross ratio quietly depends on.
\end{document}
